% GENERATED FILE. Edit canonical lessons, then run npm run generate:books.
% canonical-source-hash: fnv1a64:556203d78364053c
% canonical-lessons: GE-C120-socken, GE-C120-bluse, GE-C120-rindfleisch, GE-C120-gemuese, GE-C120-pilz

\chapter{Socks, Blouse, Beef, Vegetables, Mushroom}
\label{ch:ge-socks-blouse-beef-vegetables-mushroom}
\hlchaptermodality{120}

\begin{chapteropening}
\textbf{By the end of this chapter:} \emph{I can say socks, a blouse, beef, vegetables and a mushroom.}

Complete the last lesson of chapter 120: I can say socks, a blouse, beef, vegetables and a mushroom.
\end{chapteropening}

\section[die Socken]{die Socken — socks}

\label{lesson:GE-C120-socken}



\begin{quote}
Before the new one: say the German for a coat, then the German for shoes.
\end{quote}

\subsection*{You'll want to know: die Socken}
\textbf{die Socken} — \textquotedblleft{}socks\textquotedblright{}.

\begin{cousinweb}
\textbf{Socken} is the plural of \emph{Socke} $\leftarrow$ Old High German \emph{soc} $\leftarrow$ Latin \emph{soccus}, \textquotedblleft{}slipper\textquotedblright{} (EtymDB 2.1) — an early loan, like English \emph{sock}.
\end{cousinweb}

\begin{grammarlens}[title={What you've built}]
One more word for this chapter.
\end{grammarlens}

\subsection*{Your turn}
\begin{itemize}
\raggedright
  \item \emph{Say it:} \emph{die Socken}
  \item \emph{Say it:} \emph{die Socken}, once more
  \item \emph{Say it:} say \emph{die Socken}
  \item \emph{Recall:} say the German for a coat, then the German for shoes, then say \emph{die Socken} again
\end{itemize}

\begin{tcolorbox}[breakable,colback=teal!4,colframe=teal!35!black,title={Before you move on}]
What does die Socken mean? (\textquotedblleft{}Socks\textquotedblright{}.) Say it once more.
\end{tcolorbox}
\section[die Bluse]{die Bluse — a blouse}

\label{lesson:GE-C120-bluse}



\begin{quote}
Before the new one: say the German for shoes, then the German for socks.
\end{quote}

\subsection*{You'll want to know: die Bluse}
\textbf{die Bluse} — \textquotedblleft{}a blouse\textquotedblright{}.

\begin{cousinweb}
\textbf{Bluse} $\leftarrow$ French \emph{blouse} (EtymDB 2.1) — borrowed, like English \emph{blouse}. Where the French word came from is uncertain.
\end{cousinweb}

\begin{grammarlens}[title={What you've built}]
One more word for this chapter.
\end{grammarlens}

\subsection*{Your turn}
\begin{itemize}
\raggedright
  \item \emph{Say it:} \emph{die Bluse}
  \item \emph{Say it:} \emph{die Bluse}, once more
  \item \emph{Say it:} say \emph{die Bluse}
  \item \emph{Recall:} say the German for shoes, then the German for socks, then say \emph{die Bluse} again
\end{itemize}

\begin{tcolorbox}[breakable,colback=teal!4,colframe=teal!35!black,title={Before you move on}]
What does die Bluse mean? (\textquotedblleft{}A blouse\textquotedblright{}.) Say it once more.
\end{tcolorbox}
\section[das Rindfleisch]{das Rindfleisch — beef}

\label{lesson:GE-C120-rindfleisch}



\begin{quote}
Before the new one: say the German for socks, then the German for a blouse.
\end{quote}

\subsection*{You'll want to know: das Rindfleisch}
\textbf{das Rindfleisch} — \textquotedblleft{}beef\textquotedblright{}.

\begin{cousinweb}
\textbf{Rindfleisch} is the word for \textquotedblleft{}cattle\textquotedblright{}, Old High German \emph{hrind} $\leftarrow$ Proto-Germanic \emph{*hrinþaz}, + \emph{Fleisch}, \textquotedblleft{}meat\textquotedblright{} (EtymDB 2.1).
\end{cousinweb}

\begin{grammarlens}[title={What you've built}]
One more word for this chapter.
\end{grammarlens}

\subsection*{Your turn}
\begin{itemize}
\raggedright
  \item \emph{Say it:} \emph{das Rindfleisch}
  \item \emph{Say it:} \emph{das Rindfleisch}, once more
  \item \emph{Say it:} say \emph{das Rindfleisch}
  \item \emph{Recall:} say the German for socks, then the German for a blouse, then say \emph{das Rindfleisch} again
\end{itemize}

\begin{tcolorbox}[breakable,colback=teal!4,colframe=teal!35!black,title={Before you move on}]
What does das Rindfleisch mean? (\textquotedblleft{}Beef\textquotedblright{}.) Say it once more.
\end{tcolorbox}
\section[das Gemüse]{das Gemüse — vegetables}

\label{lesson:GE-C120-gemuese}



\begin{quote}
Before the new one: say the German for a blouse, then the German for beef.
\end{quote}

\subsection*{You'll want to know: das Gemüse}
\textbf{das Gemüse} — \textquotedblleft{}vegetables\textquotedblright{}.

\begin{cousinweb}
\textbf{Gemüse} $\leftarrow$ Middle High German \emph{gemüese}: \emph{ge-} + the word for \textquotedblleft{}mash, food\textquotedblright{}, Old High German \emph{muos} $\leftarrow$ Proto-Germanic \emph{*mōsą} (EtymDB 2.1). Vegetables were what went into the pot.
\end{cousinweb}

\begin{grammarlens}[title={What you've built}]
One more word for this chapter.
\end{grammarlens}

\subsection*{Your turn}
\begin{itemize}
\raggedright
  \item \emph{Say it:} \emph{das Gemüse}
  \item \emph{Say it:} \emph{das Gemüse}, once more
  \item \emph{Say it:} say \emph{das Gemüse}
  \item \emph{Recall:} say the German for a blouse, then the German for beef, then say \emph{das Gemüse} again
\end{itemize}

\begin{tcolorbox}[breakable,colback=teal!4,colframe=teal!35!black,title={Before you move on}]
What does das Gemüse mean? (\textquotedblleft{}Vegetables\textquotedblright{}.) Say it once more.
\end{tcolorbox}
\section[der Pilz]{der Pilz — a mushroom}

\label{lesson:GE-C120-pilz}



\begin{quote}
Before the new one: say the German for beef, then the German for vegetables.
\end{quote}

\subsection*{You'll want to know: der Pilz}
\textbf{der Pilz} — \textquotedblleft{}a mushroom\textquotedblright{}.

\begin{cousinweb}
\textbf{Pilz} $\leftarrow$ Old High German \emph{buliz}, from Latin \emph{bōlētus}, \textquotedblleft{}edible mushroom\textquotedblright{} $\leftarrow$ Greek \emph{bōlítēs} — borrowed. EtymDB 2.1 has the Latin and Greek step.
\end{cousinweb}

\begin{grammarlens}[title={What you've built}]
One more word for this chapter.
\end{grammarlens}

\subsection*{Your turn}
\begin{itemize}
\raggedright
  \item \emph{Say it:} \emph{der Pilz}
  \item \emph{Say it:} \emph{der Pilz}, once more
  \item \emph{Say it:} say \emph{der Pilz}
  \item \emph{Recall:} say the German for beef, then the German for vegetables, then say \emph{der Pilz} again
\end{itemize}

\begin{tcolorbox}[breakable,colback=teal!4,colframe=teal!35!black,title={Before you move on}]
What does der Pilz mean? (\textquotedblleft{}A mushroom\textquotedblright{}.) Say it once more.
\end{tcolorbox}
